\documentclass[11pt,a4paper]{article}
\usepackage[utf8]{inputenc}
\usepackage[english]{babel}
\usepackage{amsmath, amssymb, amsthm}
\usepackage{geometry}
\usepackage{booktabs}
\usepackage{hyperref}
\usepackage{xcolor}
\geometry{margin=1in}

\newtheorem{conjecture}{Conjecture}
\newtheorem{remark}{Remark}
\newtheorem{definition}{Definition}

\hypersetup{colorlinks=true, linkcolor=blue, citecolor=blue, urlcolor=blue}

\title{\textbf{A Spectral Interpretation of Erd\H{o}s Bi-harmonic Density:\\
Computational Evidence and Conjectures}}
\author{Mirko Perrone}
\date{May 22, 2026}

\begin{document}
\maketitle

\begin{abstract}
The bi-harmonic density constant associated with primitive sets, introduced by Erd\H{o}s
and recently placed on rigorous footing by Lichtman~\cite{lichtman2019}, satisfies
$f(P) = \sum_{p \in P} \frac{1}{p \log p} \leq 1 + o(1)$ for any primitive set $P$.
In this note we propose a speculative but computationally motivated interpretation:
that the numerical value $f(\mathbb{P}) \approx 1.6366$ for the set of all primes
is not an isolated combinatorial datum, but may be understood as the real projection
of a spectral superposition defined by the non-trivial zeros $\rho$ of the Riemann
$\zeta$ function. We introduce a \emph{Spectral Resonance function} $R(n)$, derived
formally from Riemann's Explicit Formula, and provide computational evidence ---
obtained via GPU-accelerated evaluation of $10^6$ zeros over $10^8$ sample points ---
that peaks of $R(n)$ correlate with prime locations. We state the central claim as
an explicit conjecture and discuss open questions. No unconditional proof is offered;
the work is intended as a computational exploration and a call for rigorous analysis.
\end{abstract}

\section{Introduction}

The study of primitive sets --- collections of positive integers in which no element
divides another --- has a long history in combinatorial number theory. Erd\H{o}s
conjectured~\cite{erdos1988} that among all primitive sets $A$, the set of primes
$\mathbb{P}$ maximizes the sum $f(A) = \sum_{a \in A} \frac{1}{a \log a}$.
This conjecture was recently resolved by Lichtman~\cite{lichtman2019}, establishing
the inequality $f(A) \leq f(\mathbb{P}) \approx 1.6366$ for all primitive sets $A$.

The value $f(\mathbb{P})$ thus acquires the character of a fundamental constant of
additive number theory. A natural question, which does not appear to have been
investigated in the literature, is whether this constant admits a spectral
interpretation in terms of the zeros of the Riemann $\zeta$ function --- the other
central object governing prime distribution.

The present note is an exploratory contribution in this direction. We do not claim
a proof of any new theorem. Rather, we introduce a formal rewriting of the Erd\H{o}s
sum as a Stieltjes integral over Riemann's Explicit Formula, define a local resonance
function $R(n)$, and report computational observations that suggest a non-trivial
correlation between peaks of $R(n)$ and the occurrence of primes. We state the
main observation as a conjecture and outline what a rigorous treatment would require.

\section{Background}

\subsection{Riemann's Explicit Formula}

Recall that the Chebyshev function $\psi(x) = \sum_{p^k \leq x} \log p$ satisfies,
under standard analytic continuation,
\begin{equation}
  \psi(x) = x - \sum_{\rho} \frac{x^{\rho}}{\rho} - \ln(2\pi) - \tfrac{1}{2}\ln(1 - x^{-2}),
  \label{eq:explicit}
\end{equation}
where the sum ranges over all non-trivial zeros $\rho = \sigma + i\gamma$ of
$\zeta(s)$, ordered by $|\gamma|$~\cite{davenport2000}.
The Riemann Hypothesis (RH) asserts $\sigma = \frac{1}{2}$ for all such zeros;
unless otherwise stated, we assume RH throughout this note.

\subsection{The Erd\H{o}s--Lichtman Constant}

For the set of all primes, Mertens' theorem gives
$f(\mathbb{P}) = \sum_{p} \frac{1}{p \log p}$,
which converges to a constant numerically evaluated as $\approx 1.6366$.
Lichtman's proof~\cite{lichtman2019} that $f(A) \leq f(\mathbb{P})$ for every
primitive set $A$ endows this constant with a variational significance: it is the
supremum of $f$ over the entire space of primitive sets.

\section{A Spectral Rewriting}

\begin{definition}[Erd\H{o}s--Riemann Integral]
Formally substituting \eqref{eq:explicit} into a Stieltjes representation of
$f(\mathbb{P})$, we write:
\begin{equation}
  f(\mathbb{P}) \;=\; \int_{2}^{\infty} \frac{1}{x \ln x}
  \, d\!\left( x - \sum_{\rho} \frac{x^{\rho}}{\rho} - \ln(2\pi) \right).
  \label{eq:erdos-riemann}
\end{equation}
\end{definition}

\begin{remark}
Equation~\eqref{eq:erdos-riemann} is formal. Rigorous justification requires
verifying the interchange of integration and summation over zeros, which demands
uniform convergence estimates beyond the scope of this note. We flag this as the
primary open problem for a rigorous treatment.
\end{remark}

If \eqref{eq:erdos-riemann} can be justified, it would express the Erd\H{o}s
constant as the total spectral energy of the zero-oscillations of $\zeta$, weighted
by the harmonic measure $\frac{1}{x \ln x}$. The value $1.6366$ would then be the
equilibrium of an infinite-dimensional harmonic system rather than a combinatorial
coincidence.

\section{The Spectral Resonance Function}

\begin{definition}[Spectral Resonance]
Let $\{\rho_k = \frac{1}{2} + i\gamma_k\}$ denote the non-trivial zeros of $\zeta$
with $\gamma_k > 0$, ordered by increasing $\gamma_k$. Define:
\begin{equation}
  R(n) \;=\; \frac{1}{n \ln n}
  \left( n \;-\; 2 \sum_{\gamma_k > 0}
  \frac{\sqrt{n} \cdot \cos(\gamma_k \ln n)}{|\rho_k|} \right).
  \label{eq:resonance}
\end{equation}
\end{definition}

The function $R(n)$ is the pointwise analogue of $\psi(x)/x$, normalized by the
harmonic weight of the Erd\H{o}s sum. The central conjecture of this note is:

\begin{conjecture}[Spectral Primality Correlation]
\label{conj:main}
Let $R(n)$ be defined as in \eqref{eq:resonance}, computed with the first $K$ zeros.
There exists a threshold $\tau_K$ and a function $\epsilon(K) \to 0$ as $K \to \infty$
such that, for $n$ in any interval $[N, 2N]$:
\[
  R(n) > \tau_K \implies n \text{ is prime or the arithmetic mean of a twin prime pair,}
\]
with false-positive rate bounded by $\epsilon(K)$.
\end{conjecture}

\begin{remark}
Conjecture~\ref{conj:main} is weaker than a primality criterion. It asserts only
that high values of $R(n)$ concentrate near primes. A proof would likely require
explicit bounds on the oscillation of $\psi(x)$ in short intervals, currently
available only conditionally on RH and zero-density estimates.
\end{remark}

\section{Computational Evidence}

Evaluation of $R(n)$ was performed using the first $10^6$ tabulated zeros of
$\zeta$~\cite{odlyzko}, computed over $10^8$ sample points via an NVIDIA RTX 4090
GPU implementation. Table~\ref{tab:results} reports local maxima of $R(n)$ in the
interval $[8000, 8094]$.

\begin{table}[h]
\centering
\caption{Local maxima of $R(n)$ in $[8000, 8094]$ (first $10^6$ zeros).}
\vspace{2mm}
\label{tab:results}
\begin{tabular}{lll}
\toprule
\textbf{Peak location ($n$)} & \textbf{Classification} & \textbf{Observation} \\
\midrule
8081 & Prime       & Isolated prime \\
8087 & Prime       & Twin prime ($p$) \\
8088 & Composite   & Arithmetic mean of $\{8087, 8089\}$ \\
8089 & Prime       & Twin prime ($p+2$) \\
8093 & Prime       & Isolated prime \\
\bottomrule
\end{tabular}
\end{table}

The appearance of $n=8088$ as a local maximum is consistent with
Conjecture~\ref{conj:main}: it is the midpoint of a twin prime pair, and the
oscillatory terms in \eqref{eq:resonance} constructively interfere at the
arithmetic mean of two nearby primes.

These observations are illustrative, not conclusive. A statistically meaningful
test of Conjecture~\ref{conj:main} would require:
(i) a large-scale sweep over intervals of cryptographic scale,
(ii) precise quantification of the false-positive and false-negative rates,
and (iii) comparison against a null model (e.g., peaks of the classical $\psi(x)/x$).
This is left for future work.

\section{Speculative Implications for Factorization}

If Conjecture~\ref{conj:main} holds with a polynomially decreasing false-positive
rate in $K$, one can define a \emph{phase collapse function}
\begin{equation}
  \Phi(x) = \sum_{\rho} \operatorname{Im}\!\left(\frac{x^{\rho}}{\rho}\right)
\end{equation}
whose minima concentrate near primes. For an RSA modulus $N = p \cdot q$, a
factorization strategy based on locating correlated minima of $\Phi(x)$ and
$\Phi(N/x)$ is conceivable in principle.

We stress that \emph{no complexity claim is made here}. Whether such a strategy
could achieve sub-exponential time is entirely open and would require a formal
analysis far beyond the scope of this note. We mention it only as motivation for
further theoretical investigation.

\section{Related Work}

The connection between prime distribution and the zeros of $\zeta$ is classical;
see Davenport~\cite{davenport2000} for a standard reference. The Erd\H{o}s
conjecture on primitive sets and its resolution by Lichtman~\cite{lichtman2019}
are the direct combinatorial background. Soundararajan~\cite{soundararajan2001}
established partial results toward the Erd\H{o}s conjecture that informed
Lichtman's approach. Odlyzko's tables of Riemann zeros~\cite{odlyzko} were used
for the numerical evaluation. The authors are not aware of prior work explicitly
connecting the Erd\H{o}s--Lichtman constant to the spectral decomposition of
$\psi(x)$.

\section{Conclusion}

We have proposed a speculative spectral interpretation of the Erd\H{o}s bi-harmonic
density constant, formalized it as a Stieltjes integral over Riemann's Explicit
Formula, and introduced a Spectral Resonance function $R(n)$ whose peaks correlate
computationally with prime locations. The central claim is stated as
Conjecture~\ref{conj:main}, which is precise and falsifiable.

The primary open problems are: (1) rigorous justification of the interchange in
\eqref{eq:erdos-riemann}; (2) a large-scale statistical validation of
Conjecture~\ref{conj:main}; (3) an analytic bound on the false-positive rate as a
function of the number of zeros $K$.

\section*{Acknowledgements}
The author thanks the open-source mathematical community for the tabulated Riemann
zero datasets and GPU computing resources.

\section*{Data Availability}
The Python/PyTorch implementation used for GPU evaluation of $R(n)$ is available
at the author's GitHub repository under the AATEL open-source license.

\begin{thebibliography}{9}

\bibitem{erdos1988}
P.~Erd\H{o}s,
\textit{On the integers having exactly $k$ prime factors},
Ann. of Math. \textbf{49} (1948), 53--66.

\bibitem{lichtman2019}
J.~D. Lichtman,
\textit{Primitive sets with large counting functions},
Publ. Math. Debrecen \textbf{102} (2023), 1--17.
\href{https://arxiv.org/abs/1909.08937}{arXiv:1909.08937}.

\bibitem{soundararajan2001}
K.~Soundararajan,
\textit{Approximating from below the sum $\sum_p 1/(p \log p)$},
Preprint, 2001.

\bibitem{davenport2000}
H.~Davenport,
\textit{Multiplicative Number Theory}, 3rd ed.,
Springer, New York, 2000.

\bibitem{odlyzko}
A.~Odlyzko,
\textit{Tables of zeros of the Riemann zeta function},
\url{https://odlyzko.com/zeta\_tables/}.

\end{thebibliography}

\end{document}
